\section{The smallest orbit controls uniform concentration}
\label{sec:strata}
In this section $G$ is a compact connected Lie group, $U:G\to O(V)$,
and $V^G=\{0\}$. Put
\begin{equation}\label{eq:min-orbit}
 s(U)=\min_{u\in\mathbb S(V)}\dim(G u).
\end{equation}
The minimum is over \emph{all} directions. It can be smaller than the
dimension of a generic orbit. Since $G$ is connected, $s(U)\ge1$.

\begin{theorem}[Uniform concentration across orbit types]\label{thm:uniform-orbits}
There is a finite constant $A_U\ge1$ such that
\begin{equation}\label{eq:uniform-orbits}
 \sup_{u,v\in\mathbb S(V)}m_G\{g:\norm{U_gu-v}\le r\}
       \le A_U r^{s(U)}\qquad(0<r\le1).
\end{equation}
This exponent is optimal among uniform power upper bounds: for a unit
$u_0$ with $\dim(G u_0)=s(U)$ there are $c,r_0>0$ such that
\begin{equation}\label{eq:orbit-sharpness}
 m_G\{g:\norm{U_gu_0-u_0}\le r\}\ge c r^{s(U)}
 \qquad(0<r<r_0).
\end{equation}
No regular-orbit, constant-stabilizer, or simple-spectrum hypothesis
is needed in \eqref{eq:uniform-orbits}.
\end{theorem}
\begin{proof}
Write $m=\dim G$, $s=s(U)$ and let $\mathfrak g$ be the Lie algebra.
For $u\in\mathbb S(V)$ the differential
\[
 A_u:\mathfrak g\longrightarrow V,\qquad X\longmapsto dU(X)u
\]
has rank $\dim(G u)\ge s$. Its $s$th singular value is continuous in
$u$ and is strictly positive on the compact sphere. Its minimum is
therefore positive. We give a chart argument that also records why
variation of the stabilizer causes no loss in the exponent.

Fix $(g_0,u_0)\in G\times\mathbb S(V)$. In a smooth coordinate chart
$g=\psi(z,w)$, where $z\in\R^s$ and $w\in\R^{m-s}$, select $s$
coordinate directions and a linear projection $L:V\to\R^s$ for which
\[
 D_z\bigl[L U_{\psi(z,w)}u\bigr]_{(g_0,u_0)}=M
\]
is invertible. This is possible by the rank bound. By shrinking to a
product of bounded open rectangles $Z\times W$ and a neighborhood
$J$ of $u_0$, arrange
\[
 \norm{D_z[L U_{\psi(z,w)}u]-M}_{\rm op}
          \le\tfrac12\sigma_{\min}(M)
       \quad (z,w,u)\in Z\times W\times J.
\]
Take $Z$ convex. For fixed $(w,u)$, integration along the segment from
$z$ to $z'$ gives
\begin{equation}\label{eq:chart-co-lipschitz}
 \norm{L U_{\psi(z,w)}u-L U_{\psi(z',w)}u}
       \ge\tfrac12\sigma_{\min}(M)\norm{z-z'}.
\end{equation}
If both points lie in the inverse image of the radius-$r$ ball around
$v$, their projected distance is at most $2\norm L r$. This inverse
image, on the fixed $w$ slice, consequently has diameter at most
$4\norm L r/\sigma_{\min}(M)$, and hence $s$-dimensional Lebesgue
measure at most a constant times $r^s$. Haar measure has a bounded
smooth density in the chosen relatively compact chart. Integration
in the remaining $m-s$ coordinates yields the same bound there.

Choose finitely many of these product neighborhoods covering
$G\times\mathbb S(V)$. For each fixed $u$, those whose $J$ contains
$u$ cover $G$; summing the preceding bounds proves
\eqref{eq:uniform-orbits}. All constants are independent of the
center $v$. An empty inverse image contributes zero. This use of a
finite cover, rather than a separate density estimate on every orbit,
is what controls neighborhoods of singular orbit types.

For the converse, $G u_0\simeq G/G_{u_0}$ is a compact smooth embedded
$s$-dimensional submanifold. The pushforward of Haar probability is
its normalized invariant Riemannian volume: the induced metric is
$G$-invariant and the invariant probability on this homogeneous space
is unique. A smooth local graph at $u_0$ has volume comparable to
$r^s$ in a small ambient ball. This proves \eqref{eq:orbit-sharpness}.
If a uniform estimate with exponent $s'>s$ held, comparison at $u_0$
as $r\downarrow0$ would contradict it.
\end{proof}

The rank computation in \eqref{eq:min-orbit} is intrinsic to the
representation. For example, given algebraic skew matrices
$B_1,\ldots,B_m$ spanning its infinitesimal action, the condition
$s(U)\le r$ is the existence of a real vector $u$ with
$\sum u_i^2=1$ and all $(r+1)$-minors of $[B_1u\ \cdots\ B_mu]$
zero. Thus it is an explicit real-algebraic decision problem for
fully specified algebraic data; quantifier elimination applies
\cite{BPR}. This is not a test for a spectral gap, and no efficient
algorithm in a succinct representation is asserted.

\begin{corollary}[Intrinsic occupation exponent]\label{cor:intrinsic-budget}
Suppose the representation satisfies \eqref{eq:gap}. In the setting
of Lemma~\ref{lem:terminal}, $q_N\ge\kappa>0$ implies
\begin{equation}\label{eq:intrinsic-budget}
 \sum_{\substack{0\le t<N\\K_t\le k}} g_t
       \le 2048D(16A_U)^{2/s(U)}
                    \kappa^{-3-2/s(U)}k^{2/s(U)}.
\end{equation}
The estimate allows arbitrary hidden continuations and imposes no
bound on the other cuts.
\end{corollary}
\begin{proof}
Apply Theorem~\ref{thm:occupation} with the uniform estimate
\eqref{eq:uniform-orbits}. In particular the conditional means may
change orbit type at every step; their directions still satisfy the
same estimate.
\end{proof}
